%!TEX TS-program = pdflatex

\documentclass{article}
%% Version flag for shared include files

\usepackage{amsmath}

\title{Lorem ipsum dolor sit amet}
\author{Ann Author}

\begin{document}

\maketitle

\begin{abstract}
Lorem ipsum dolor sit amet, consectetuer adipiscing elit. Aenean commodo ligula
eget dolor. Aenean massa. Cum sociis natoque penatibus et magnis dis parturient
montes, nascetur ridiculus mus. Donec quam felis, ultricies nec, pellentesque
eu, pretium quis, sem. Nulla consequat massa quis enim. Donec pede justo,
fringilla vel, aliquet nec, vulputate eget, arcu. In enim justo, rhoncus ut,
imperdiet a, venenatis vitae, justo. Nullam dictum felis eu pede mollis pretium.
\end{abstract}


\section{Introduction}
\label{sec:introduction}
Lorem ipsum dolor sit amet, consectetuer adipiscing elit. Aenean commodo ligula
eget dolor. Aenean massa. Cum sociis natoque penatibus et magnis dis parturient
montes, nascetur ridiculus mus. Donec quam felis, ultricies nec, pellentesque
eu, pretium quis, sem. Nulla consequat massa quis enim. Donec pede justo,
fringilla vel, aliquet nec, vulputate eget, arcu. In enim justo, rhoncus ut,
imperdiet a, venenatis vitae, justo. Nullam dictum felis eu pede mollis pretium.

\section{Definition of HCENCone}
\label{sec:definition}

Lorem ipsum dolor sit amet, consectetuer adipiscing elit. Aenean commodo ligula
eget dolor. Aenean massa. Cum sociis natoque penatibus et magnis dis parturient
montes, nascetur ridiculus mus. Donec quam felis, ultricies nec, pellentesque
eu, pretium quis, sem. Nulla consequat massa quis enim. Donec pede justo,
fringilla vel, aliquet nec, vulputate eget, arcu. In enim justo, rhoncus ut,
imperdiet a, venenatis vitae, justo. Nullam dictum felis eu pede mollis pretium.

\paragraph{The Key Derivation.}
The master key is a pair \(K_M=(K_a,K_b)\) of blocks.
First, a block \(L\) is derived from \(K_M\), and from it and \(K_a\) six further blocks, as follows:
\begin{equation}
\label{eq:key-derivation}
 L=E_{K_a}(K_b)\oplus K_a\enspace,\qquad
 d_j=E_{L\oplus\mathrm{bin}_n(j)}(K_b \oplus K_a)\enspace,\qquad 0\leq j\leq5\enspace.
\end{equation}
The six derived blocks have three uses:
\begin{itemize}
\item \(h_1 := d_0\) and \(h_2 := d_1\) are the hash keys;
\item \((K_1,K_2,K_3) := (d_2,d_3,d_4)\) are the base keys of the head construction;
\item \(K_S := d_5\) is the base key of the stream.
\end{itemize}
Lorem ipsum dolor sit amet, consectetuer adipiscing elit. Aenean commodo ligula
eget dolor. Aenean massa. Cum sociis natoque penatibus et magnis dis parturient
montes, nascetur ridiculus mus. Donec quam felis, ultricies nec, pellentesque
eu, pretium quis, sem. Nulla consequat massa quis enim. Donec pede justo,
fringilla vel, aliquet nec, vulputate eget, arcu. In enim justo, rhoncus ut,
imperdiet a, venenatis vitae, justo. Nullam dictum felis eu pede mollis pretium.


\newenvironment{colorthis}[1]{\color{#1}}{\ignorespacesafterend}

\paragraph{Alternative Key Derivations.}
The master key is a pair \(K_M=(K_a,K_b)\) of blocks.
First, a block \(L\) is derived from \(K_M\), and from it and \(K_a\) six further blocks, as follows:
\begin{align}
\label{eq:key-derivation-2}
L=E_{K_a}(K_b)\enspace,\qquad
 d_j&=E_{L\oplus\mathrm{bin}_n(j)}(K_a)\oplus K_b\enspace,\qquad 0\leq j\leq5\enspace. \\
 \label{eq:key-derivation-3}
 L=E_{K_a}(K_b)\enspace,\qquad
 d_j&=E_{L}(K_a \oplus\mathrm{bin}_n(j))\oplus K_b\enspace,\qquad 0\leq j\leq5\enspace.
\end{align}

\paragraph{The Hash.}
Lorem ipsum dolor sit amet, consectetuer adipiscing elit. Aenean commodo ligula
eget dolor. Aenean massa. Cum sociis natoque penatibus et magnis dis parturient
montes, nascetur ridiculus mus. Donec quam felis, ultricies nec, pellentesque
eu, pretium quis, sem. Nulla consequat massa quis enim. Donec pede justo,
fringilla vel, aliquet nec, vulputate eget, arcu. In enim justo, rhoncus ut,
imperdiet a, venenatis vitae, justo. Nullam dictum felis eu pede mollis pretium.
\[
 \textrm{MHASH}_n(h,B_1\Vert\cdots\Vert B_r)=\sum_{i=1}^rB_i\eta^{r-i+1}\enspace,
\]
Lorem ipsum dolor sit amet, consectetuer adipiscing elit. Aenean commodo ligula
eget dolor. Aenean massa. Cum sociis natoque penatibus et magnis dis parturient
montes, nascetur ridiculus mus. Donec quam felis, ultricies nec, pellentesque
eu, pretium quis, sem. Nulla consequat massa quis enim. Donec pede justo,
fringilla vel, aliquet nec, vulputate eget, arcu. In enim justo, rhoncus ut,
imperdiet a, venenatis vitae, justo. Nullam dictum felis eu pede mollis pretium.
The mode evaluates this hash under both hash keys and uses the two values separately:
\[
 H^s(T,U)=\bigl(H^s_{h_1}(T,U),\,H^s_{h_2}(T,U)\bigr)\in\{0,1\}^{2n}\enspace.
\]


\end{document}
